\subsection{映射的逆系统}

命题1。设I是一个有序集，令 \((\mathbf{E}_{\alpha}, f_{\alpha \beta})\) 是相对于 I 的集合逆系统，令 \(\mathbf{E} = \lim \mathbf{E}_{\alpha}\) 为其逆极限，且对每个 \(\alpha \in I\) 令

\[
f _ {\alpha}: \mathbf {E} \rightarrow \mathbf {E} _ {\alpha}
\]

为典范映射。对每个 \(\alpha \in \mathbf{I}\) ，令 \(u_{\alpha}\) 为从一个集合 \(\mathbf{F}\) 到 \(\mathbf{E}_{\alpha}\) 中的一个映射，使得

\[
f _ {\alpha \beta} \circ u _ {\beta} = u _ {x}\tag{5}
\]

每当 \(\alpha \leqslant \beta\) .

于是：

\begin{enumerate}
\def\labelenumi{(\alph{enumi})}
\tightlist
\item
  存在唯一的映射 \(u\) （从 \(\mathbf{F}\) 到 \(\mathbf{E}\) 中），使得
\end{enumerate}

\[
u _ {\alpha} = f _ {\alpha} \circ u\tag{6}
\]

\[
\text { 对   所有 } \alpha \in \mathbf {I};
\]

\begin{enumerate}
\def\labelenumi{(\alph{enumi})}
\setcounter{enumi}{1}
\tightlist
\item
  映射 \(u\) 是单射的，当且仅当，对于每一对不同的元素 \(y, z\) （\(\mathbf{F}\) 的元素）, 存在 \(\alpha \in \mathbf{I}\) 使得 \(u_{\alpha}(y) \neq u_{\alpha}(z)\) .
\end{enumerate}

因为关系 \(u_{\alpha} = f_{\alpha} \circ u\) 意味着，对于每个 \(y \in F\) 我们有

\[
\operatorname{pr} _ {\alpha} (u (\gamma)) = u _ {\alpha} (\gamma);
\]

该元素 \(u(y) \in \prod_{\alpha \in I} E_{\alpha}\) 由 \(u(y) = (u_{\alpha}(y))_{\alpha \in I}\) 唯一确定。仍需证明 \(u(y) \in E\) （对于所有 \(y \in F\) ），换言之，即

\[
\operatorname{pr} _ {\alpha} (u (y)) = f _ {\alpha \beta} (\operatorname{pr} _ {\beta} (u (y)))
\]

每当 \(\alpha \leqslant \beta\) 。但这可以写成形式

\[
u _ {\alpha} (y) = f _ {\alpha \beta} (u _ {\beta} (y))
\]

因此由(5)得出。命题的第二部分直接由定义得出。

推论1. 设 \((\mathbf{E}_{\alpha}, f_{\alpha \beta})\) 和 \((\mathbf{F}_{\alpha}, g_{\alpha \beta})\) 是相对于同一指标集 \(\mathbf{I}\) 的两个集合逆系统；令 \(\mathbf{E} = \lim \mathbf{E}_{\alpha}\) , \(\mathbf{F} = \lim \mathbf{F}_{\alpha}\) ，并且令 \(f_{\alpha}\) （相应地， \(g_{\alpha}\) ）为从 \(\mathbf{E}\) 到 \(\overleftarrow{\mathbf{E}_{\alpha}}\) 中（相应地，从 \(\overleftarrow{\mathbf{F}}\) 到 \(\mathbf{F}_{\alpha}\) 中）的典范映射，对于每个 \(\alpha \in \mathbf{I}\) 。对于每个 \(\alpha \in \mathbf{I}\) ，令 \(u_{\alpha}\) 为从 \(\mathbf{E}\) 到 \(\mathbf{F}_{\alpha}\) 中的一个映射，使得图表

\[
\begin{array}{c} \mathbf {E} _ {\beta} \xrightarrow {u _ {\beta}} \mathbf {F} _ {\beta} \\ f _ {\alpha \beta} \Big \downarrow \qquad \qquad \qquad \qquad \Big \downarrow g _ {\alpha \beta} \\ \mathbf {E} _ {\alpha} \xrightarrow [ u _ {\alpha} ]{} \mathbf {F} _ {\alpha} \end{array}
\]

是交换的 (*) （每当 \(\alpha \leqslant \beta\) ）。则存在唯一的映射 \(u: \mathbf{E} \to \mathbf{F}\) ，使得对于每一个 \(\alpha \in \mathbf{I}\) ，图表

\[
\begin{array}{c} \mathbf {E} \xrightarrow {u} \mathbf {F} \\ f _ {\alpha} \Big \downarrow \qquad \qquad \qquad \qquad \qquad \qquad \qquad \qquad \qquad \Big \downarrow g _ {\alpha} \\ \mathbf {E} _ {\alpha} \xrightarrow [ u _ {\alpha} ]{} \mathbf {F} _ {\alpha} \end{array}
\]

是交换的。

设 \(v_{\alpha} = u_{\alpha} \circ f_{\alpha}\) 。如果 \(\alpha \leqslant \beta\) ，由 (2) 我们有

\[
g _ {\alpha \beta} \circ v _ {\beta} = g _ {\alpha \beta} \circ u _ {\beta} \circ f _ {\beta} = u _ {\alpha} \circ f _ {\alpha \beta} \circ f _ {\beta} = u _ {\alpha} \circ f _ {\alpha} = v _ {\alpha},
\]

因此我们可以将命题 1 应用于映射 \(v_{\alpha}\) ；因此

一个映射 \(u: E \rightarrow F\) 的存在性和唯一性，使得

\[
g _ {\alpha} \circ u = v _ {\alpha} = u _ {\alpha} \circ f _ {\alpha}
\]

对于每一个 \(\alpha \in I\) .

一族映射 \(u_{\alpha}: \mathbf{E}_{\alpha} \to \mathbf{F}_{\alpha}\) ，若满足推论1的条件，则称为从 \((\mathbf{E}_{\alpha}, f_{\alpha\beta})\) 到 \((\mathbf{F}_{\alpha}, g_{\alpha\beta})\) 中的映射逆系统。在推论1中定义的映射 \(u\) 称为族 \((u_{\alpha})\) 的逆极限，并记作 \(u = \lim_{\leftarrow} u_{\alpha}\) （当不致引起混淆时）。

推论2。设 \((\mathbf{E}_{\alpha}, f_{\alpha \beta})\) , \((\mathbf{F}_{\alpha}, g_{\alpha \beta})\) , \((\mathbf{G}_{\alpha}, h_{\alpha \beta})\) 是相对于同一个指标集I的三个集合逆系统；设 \(\mathbf{E} = \varprojlim \mathbf{E}_{\alpha}\) , \(\mathbf{F} = \varprojlim \mathbf{F}_{\alpha}\) , \(\mathbf{G} = \varprojlim \mathbf{G}_{\alpha}\) ，并且令 \(f_{\alpha}\) （相应地， \(g_{\alpha}, h_{\alpha}\) ）为从 \(\mathbf{E}\) （相应地， \(\mathbf{F}, \mathbf{G}\) ）到 \(\mathbf{E}_{\alpha}\) （相应地， \(\mathbf{F}_{\alpha}, \mathbf{G}_{\alpha}\) ）中的典范映射。如果 \((u_{\alpha})\) 和 \((v_{\alpha})\) 是两个映射逆系统， \(u_{\alpha}: \mathbf{E}_{\alpha} \to \mathbf{F}_{\alpha}\) , \(v_{\alpha}: \mathbf{F}_{\alpha} \to \mathbf{G}_{\alpha}\) ，则复合映射 \(v_{\alpha} \circ u_{\alpha}: \mathbf{E}_{\alpha} \to \mathbf{G}_{\alpha}\) 构成一个映射逆系统，并且我们有

\[
\varprojlim (v _ {\alpha} \circ u _ {\alpha}) = (\varprojlim v _ {\alpha}) \circ (\varprojlim u _ {\alpha}).\tag{7}
\]

因为如果我们令 \(w_{\alpha} = v_{\alpha} \circ u_{\alpha}\) ，那么如果 \(\alpha \leqslant \beta\) 我们有

\[
w _ {\alpha} \circ f _ {\alpha \beta} = v _ {\alpha} \circ (u _ {\alpha} \circ f _ {\alpha \beta}) = v _ {\alpha} \circ (g _ {\alpha \beta} \circ u _ {\beta}) = (h _ {\alpha \beta} \circ v _ {\beta}) \circ u _ {\beta} = h _ {\alpha \beta} \circ w _ {\beta},
\]

这表明 \((w_{\alpha})\) 是一个映射逆系统。此外，如果 \(u = \lim u_{\alpha}\) 并且 \(v = \lim v_{\alpha}\) ，那么 \(h_{\alpha} \circ (v \circ u) = (v_{\alpha} \circ g_{\alpha}) \circ u = (v_{\alpha} \circ u_{\alpha}) \circ f_{\alpha}\) （对于每一个 \(\alpha \in I\) ），因此，由逆极限的唯一性，我们有 \(v \circ u = \lim w_{\alpha}\) 。

设 \((\mathbf{E}_{\alpha}, f_{\alpha \beta})\) 是一个集合逆系统，并且对每个 \(\alpha \in I\) ，令 \(\mathbf{M}_{\alpha}\) 为 \(\mathbf{E}_{\alpha}\) 的一个子集。如果 \(f_{\alpha \beta}(\mathbf{M}_{\beta}) \subset \mathbf{M}_{\alpha}\) （每当 \(\alpha \leqslant \beta\) ），就称各 \(\mathbf{M}_{\alpha}\) 构成各 \(\mathbf{E}_{\alpha}\) 的一个子集逆系统。设 \(g_{\alpha \beta}\) 为从 \(\mathbf{M}_{\beta}\) 到 \(\mathbf{M}_{\alpha}\) 的映射（其中 \(\alpha \leqslant \beta\) ），其图与 \(f_{\alpha \beta}\) 在 \(\mathbf{M}_{\beta}\) 上的限制之图相同。那么显然， \((\mathbf{M}_{\alpha}, g_{\alpha \beta})\) 是一个集合逆系统，并且

\[
\varprojlim \mathbf {M} _ {\alpha} = (\varprojlim \mathbf {E} _ {\alpha}) \cap \prod_ {\alpha \in \mathbf {I}} \mathbf {M} _ {\alpha}.\tag{8}
\]

命题2. 设 \((\mathbf{E}_{\alpha}, f_{\alpha \beta})\) 和 \((\mathbf{E}_{\alpha}^{\prime}, f_{\alpha \beta}^{\prime})\) 是相对于 I 的两个集合逆系统，并令 \(u_{\alpha}\) 为从 \(\mathbf{E}_{\alpha}\) 到 \(\mathbf{E}_{\alpha}^{\prime}\) 中的一个映射（对于每一个 \(\alpha \in \mathbf{I}\) ），使得各 \(u_{\alpha}\) 构成一个映射逆系统。设 \(u = \varprojlim u_{\alpha}\) 。那么对每一个 \(x^{\prime} = (x_{\alpha}^{\prime}) \in \mathbf{E}^{\prime} = \varprojlim \mathbf{E}_{\alpha}^{\prime}\) ，诸 \(-\overleftarrow{u}_{\alpha}(x_{\alpha}^{\prime})\) 构成诸 \(\mathbf{E}_{\alpha}\) 的子集的一个逆系统，并且 \(-\overleftarrow{u}(x^{\prime}) = \varprojlim -\overleftarrow{u}_{\alpha}(x_{\alpha}^{\prime})\) .

因为如果 \(\alpha \leqslant \beta\) 并且 \(x_{\beta} \in \overline{u}_{\beta}^{1}(x_{\beta}')\) ，我们有

\[
u _ {\alpha} (f _ {\alpha \beta} (x _ {\beta})) = f _ {\alpha \beta} ^ {\prime} (u _ {\beta} (x _ {\beta})) = f _ {\alpha \beta} ^ {\prime} (x _ {\beta} ^ {\prime}) = x _ {\alpha} ^ {\prime},
\]

由此得出第一个断言；而说 \(x = (x_{\alpha}) \in \mathbf{E} = \lim \mathbf{E}_{\alpha}\) 满足 \(u(x) = x'\) 意味着，根据定义， \(u_{\alpha}(x_{\alpha}) = x_{\alpha}'\) （对于每一个 \(\alpha \in I\) ）.

推论. 如果 \(u_{\alpha}\) 对每个 \(\alpha \in I\) 都是单射的（相应地，双射的），那么 \(u\) 是单射的（相应地，双射的）。

利用命题2的记号，诸像 \(u_{\alpha}(\mathbf{E}_{\alpha})\) 也构成诸 \(\mathbf{E}_{\alpha}'\) 的子集的一个逆系统，并且我们有

\[
u (\mathrm{E}) \subset \varprojlim u _ {\alpha} (\mathrm{E} _ {\alpha});\tag{9}
\]

但这个关系的两边不一定相等（习题4）。

命题3。设I是一个预序集，令 \((\mathbf{E}_{\alpha}, f_{\alpha \beta})\) 是相对于 I 的集合逆系统，并且设 \(\mathbf{E} = \lim \mathbf{E}_{\alpha}\) 。设 J 是 I 的一个共尾子集，且 J 是右有向的，并设 \(\mathbf{E}'\) 为由 \((\mathbf{E}_{\alpha}, f_{\alpha \beta})\) 通过将指标集限制为J而得到的集合逆系统的逆极限。则典范映射 \(g\) （把E映到 \(\mathbf{E}'\) 中；第1号，公式(3)）是双射的。

对于每一个 \(\alpha \in J\) ，令 \(f_{\alpha}'\) 为典范映射 \(\mathbf{E}' \to \mathbf{E}_{\alpha}\) 。那么，由(2)和(5)， \(g\) 是唯一的从 \(\mathbf{E}\) 到 \(\mathbf{E}'\) 中的映射，使得 \(f_{\alpha} = f_{\alpha}' \circ g\) （对于所有 \(\alpha \in J\) ）(命题1)。我们将使用命题1的判据来证明 \(g\) 是单射的。如果 \(x, y\) 是 \(\mathbf{E}\) 的不同元素，那么根据定义存在 \(\alpha \in I\) 使得 \(f_{\alpha}(x) \neq f_{\alpha}(y)\) ；由于 \(J\) 在 \(I\) 中共尾，存在 \(\lambda \in J\) 使得 \(\alpha \leqslant \lambda\) ；由于 \(f_{\alpha\lambda}(f_{\lambda}(x)) \neq f_{\alpha\lambda}(f_{\lambda}(y))\) ，我们有 \(f_{\lambda}(x) \neq f_{\lambda}(y)\) 。仍需证明 \(g\) 是满射。设 \(x' = (x_{\lambda}')_{\lambda \in J}\) 为 \(\mathbf{E}'\) 的一个元素。对于每一个 \(\alpha \in I\) 存在 \(\lambda \in J\) 使得 \(\alpha \leqslant \lambda\) ，且元素 \(f_{\alpha\lambda}(x_{\lambda}')\) 不依赖于指标 \(\lambda \in J\) （满足 \(\alpha \leqslant \lambda\) ）；因为如果 \(\mu \in J\) 满足 \(\alpha \leqslant \mu\) , 存在 \(\nu \in J\) 使得 \(\lambda \leqslant \nu\) 并且 \(\mu \leqslant \nu\) ，因此 \(f_{\alpha\lambda}(x_{\lambda}') = f_{\alpha\lambda}(f_{\lambda\nu}(x_{\nu}') = f_{\alpha\nu}(x_{\nu}')\) ，并且类似地 \(f_{\alpha\mu}(x_{\mu}') = f_{\alpha\nu}(x_{\nu}')\) . 设 \(x_{\alpha}\) 为诸 \(f_{\alpha\lambda}(x_{\lambda}')\) （对于 \(\lambda \in J\) ，满足 \(\alpha \leqslant \lambda\) ）的公共值，并且令 \(x = (x_{\alpha})_{\alpha \in I}\) 。则 \(x \in E\) ，因为如果 \(\alpha \leqslant \beta\) 并且如果 \(\lambda \in J\) 满足 \(\beta \leqslant \lambda\) ，我们有

\[
f _ {\alpha \beta} (x _ {\beta}) = f _ {\alpha \beta} (f _ {\beta \lambda} (x _ {\lambda} ^ {\prime})) = f _ {\alpha \lambda} (x _ {\lambda} ^ {\prime}) = x _ {\alpha}.
\]

最后，我们有 \(x_{\lambda}' = f_{\lambda \lambda}(x_{\lambda}')\) （对于所有 \(\lambda \in \mathbf{J}\) ），因此 \(x_{\lambda} = x_{\lambda}'\) （对于所有 \(\lambda \in \mathbf{J}\) ）；换句话说， \(f(x_{\lambda}) = x_{\lambda}',\) 使得 \(g(x) = x'\) 。因此 \(g\) 是满射，证明完成。

特别地，若 I 有最大元 \(\omega\) ，我们可以取 \(J = \{\omega\}\) ，因此 \(\lim E_{\alpha}\) 被典范地等同于 \(E_{\omega}\) .

注记

\begin{enumerate}
\def\labelenumi{(\arabic{enumi})}
\tightlist
\item
  对每个 \(\alpha \in I\) ，令 \(E_{\alpha}' = f_{\alpha}(E)\) 。由(2)，各集合 \(E_{\alpha}'\) 构成各 \(E_{\alpha}\) 的一个子集逆系统，并且显然 \(\lim E_{\alpha}' = E = \lim E_{\alpha}\) . 映射 \(f_{\alpha\beta}' : E_{\beta}' \to E_{\alpha}'\) （其中 \(\alpha \leqslant \beta\) ）的图与 \(f_{\alpha\beta}\) 在 \(E_{\beta}\) 上的限制之图相同；该映射是满射，并且有
\end{enumerate}

\[
\mathbf {E} _ {\alpha} ^ {\prime} = f _ {\alpha} (\mathbf {E}) \subset \bigcap_ {\beta \geqslant \alpha} f _ {\alpha \beta} (\mathbf {E} _ {\beta})\tag{10}
\]

对于所有 \(\alpha \in I\) .

\begin{enumerate}
\def\labelenumi{(\arabic{enumi})}
\setcounter{enumi}{1}
\item
  设 I 是一个（右）有向有序集，令 \((\mathbf{E}_{\alpha}, f_{\alpha\beta})\) 为I上的一个集合逆系统，且对每个 \(\alpha \in I\) 令 \(u_{\alpha}: F \to E_{\alpha}\) 是一个映射，使得该族 \((u_{\alpha})\) 满足公式（5）。考虑由 I 索引的逆系统 \((\mathbf{F}_{\alpha}, i_{\alpha\beta})\) ，其中 \(F_{\alpha} = F\) 对于所有 \(\alpha \in I\) 并且 \(i_{\alpha\beta}\) 是F的恒等映射。那么（第1号，例2）F被典范地等同于 \(\lim F_{\alpha}\) 。如果我们把 \(u_{\alpha}\) 视为从 \(F_{\alpha}\) 到 \(E_{\alpha}\) 中的映射，那么 \((u_{\alpha})\) 是一个映射逆系统，且由(6)定义的映射 \(u: F \to E\) 与这一映射系统的逆极限等同。因此，通过滥用语言，我们写作 \(u = \lim u_{\alpha}\) 。
\item
  令 I 为一个有序集，且令 \((\mathbf{E}_{\alpha}, f_{\alpha \beta})\) 为 I 上的一个集合逆系统。对 \(I\) 的每个有限子集 \(J\) ，让 \(F_{J}\) 为由 \((\mathbf{E}_{\alpha}, f_{\alpha \beta})\) 通过将指标集限制为J而得到的（有限）逆系统的逆极限。若J和K是I的任意两个有限子集，且满足 \(J \subset K\) ，让 \(g_{JK}\) 表示从 \(F_{K}\) 到 \(F_{J}\) 中的典范映射（3）。于是关系式(4)表明 \((F_{J}, g_{JK})\) 是相对于I的有限子集所成的有向集（关于关系 \(c\) ） \(\mathfrak{F}(I)\) 的一个集合逆系统。接下来，对于每个 \(J \in \mathfrak{F}(I)\) 让 \(h_{J}: E \to F_{J}\) 为典范映射（3）。根据（4）并借助注记（2）中提到的语言滥用， \((h_{J})\) 构成一个映射逆系统。令 \(h = \lim_{\leftarrow} h_{J}: E \to F = \lim_{\leftarrow} F_{J}\) ，且让我们证明 \(h\) 是一个双射（称为典范的）。事实上，设 \(y = (y_{J}) \in F\) 。根据定义，我们有 \(y_{J} = (x_{\alpha, J})_{\alpha \in J}\) ，其中 \(x_{\alpha, J} \in E_{\alpha}\) 对于所有 \(\alpha \in J\) 。如果 \(J \subset K\) ，则根据映射 \(g_{JK}\) 的定义，并且因为 \(y_{J} = g_{JK}(y_{K})\) ，我们有 \(x_{\alpha, J} = x_{\alpha, K}\) 对于所有 \(\alpha \in J\) 。因此，给定 \(\alpha \in I\) ，存在唯一元素 \(x_{\alpha} \in E_{\alpha}\) 使得 \(x_{\alpha} = x_{\alpha, J}\) 对于I的所有包含 \(\alpha\) 的有限子集J。如果 \(\alpha \leqslant \beta\) ，存在I的有限子集J，它同时包含 \(\alpha\) 和 \(\beta\) ；因此 \(x_{\alpha} = f_{\alpha \beta}(x_{\beta})\) 根据定义。从而 \(x = (x_{\alpha})\) 是E中满足 \(h(x) = y\) 的唯一元素。
\end{enumerate}

